\section{Method}

We formulate LLM attribution as a target-conditioned subset selection problem.
Given an input text, the method first partitions the text into contiguous chunks,
then defines a set function over arbitrary chunk subsets, and finally searches
for a compact subset whose selected text is sufficient to preserve the target
prediction while avoiding redundant evidence.  Unlike token-wise attribution
methods that score each unit independently, our method evaluates each candidate
chunk in the context of the chunks that have already been selected.  The final
explanation is therefore a small evidence set together with an importance
ranking induced by conditional marginal gains.

\subsection{Preliminaries}

Let $x$ denote an input text and $y$ denote the target label to be explained.
For classification-style LLM evaluation, each class $c \in \mathcal{Y}$ is
associated with a verbalizer string $v_c$.  We query a decoder-only LLM with the
prompt template
\[
    \mathrm{Text:}\ x\ \mathrm{Label:}
\]
and compute the conditional log-likelihood of each verbalizer.  If
$v_c=(w_1,\ldots,w_m)$ contains $m$ tokens, its normalized label score is
\[
    \ell_c(z)=\frac{1}{m}\sum_{r=1}^{m}
    \log p_{\theta}(w_r \mid \mathrm{Text:}\ z\ \mathrm{Label:}, w_{<r}),
\]
where $z$ is any perturbed or selected text.  We convert the class scores into a
label distribution by normalizing across verbalizers:
\[
    p_{\theta}(c \mid z)=
    \frac{\exp(\ell_c(z))}
    {\sum_{c' \in \mathcal{Y}}\exp(\ell_{c'}(z))}.
\]

We also use the LLM as a feature extractor.  For any text $z$, let
$\phi(z)$ be the normalized embedding obtained by taking hidden states from an
intermediate transformer layer, mean-pooling over non-padding tokens, and
$\ell_2$-normalizing the resulting vector.  In our implementation the layer is
chosen at a fixed ratio of the model depth, with ratio $0.7$ by default.  For
two texts $u$ and $v$, semantic similarity and distance are defined as
\[
    \mathrm{sim}(u,v)=
    \frac{\phi(u)\cdot \phi(v)}
    {\|\phi(u)\|_2\|\phi(v)\|_2},
    \qquad
    d(u,v)=1-\mathrm{sim}(u,v).
\]
This representation is used only through forward passes; no gradients or model
training are required by our attribution objective.

After chunking, the input is represented as an ordered partition
$\mathcal{C}=\{c_1,\ldots,c_n\}$.  Each chunk $c_i$ stores a character span
starting at $a_i$ and ending before $b_i$, with text $x[a_i:b_i]$.  The spans
are required to be contiguous, non-overlapping, and exhaustive, so concatenating
all chunks in order exactly reconstructs the original text.  For a subset
$S \subseteq \{1,\ldots,n\}$, we write $x_S$ for the text obtained by
concatenating selected chunks in their original order, and $x_{\bar{S}}$ for the
complementary text.  Empty selected or complementary text is represented by a
fixed placeholder token, so all objective terms remain well-defined.

\subsection{Adaptive Chunking}

The granularity of the chunk partition is central to text attribution.  If the
chunks are too coarse, a selected evidence unit may contain several unrelated
facts; if the chunks are too fine, local syntax is broken and the search becomes
noisy and expensive.  We therefore use a structure-aware adaptive chunker that
selects different boundary rules according to the length and structure of the
input text while maintaining exact character-level coverage.

The chunker first computes lightweight text statistics: word count, punctuation
count, newline count, and an estimated number of sentence-ending boundaries.
Given profile thresholds $\tau_s,\tau_m,\tau_l$, the input is assigned to one of
four length buckets:
\[
    \mathrm{bucket}(x)=
    \begin{cases}
    \mathrm{short}, & |x|_{\mathrm{word}} \leq \tau_s,\\
    \mathrm{medium}, & \tau_s < |x|_{\mathrm{word}} \leq \tau_m,\\
    \mathrm{long}, & \tau_m < |x|_{\mathrm{word}} \leq \tau_l,\\
    \mathrm{very\_long}, & |x|_{\mathrm{word}} > \tau_l.
    \end{cases}
\]
The default balanced profile uses increasingly permissive thresholds for short,
medium, and long texts, and the profile can be tuned without changing the
attribution objective.

For each bucket, the initial segmentation uses the strongest available
structure at the appropriate granularity.  Short texts are split by clause-like
boundaries such as commas, semicolons, colons, and dashes, because a short review
or sentence often contains several sentiment-bearing clauses.  Medium texts are
split into sentences using a sentence boundary detector when available, with a
regular-expression sentence splitter as a fallback.  Long texts are first
separated by paragraph boundaries and then split into sentences inside each
paragraph.  Very long texts start from paragraph-level spans to avoid producing
an excessively fragmented search space.

The raw spans are then repaired by a deterministic post-processing pipeline.
First, the span list is normalized so that it covers the full character range
from $0$ to $|x|$ without gaps or overlaps.  Second, invalid fragments are
merged into neighboring chunks, including whitespace-only fragments and
punctuation-only orphan chunks.  Third, very short spans are merged with
adjacent text when they fall below a minimum word or character threshold; this
prevents syntactic markers from becoming standalone candidates.  Fourth,
overlong spans are split recursively, preferring sentence boundaries, then
clause boundaries, and finally fixed-size word windows.  The maximum span length
depends on the length bucket, so short and medium texts are kept relatively
fine-grained while long documents are allowed larger coherent chunks.

Two additional guards stabilize the extremes.  For short inputs that contain
enough words but would otherwise yield too few chunks, an effective-chunk floor
splits the largest spans until a minimum number of candidates is reached.  This
avoids the degenerate case where top-percent evaluation or greedy search has no
meaningful choice.  For very long inputs, a fragmentation guard merges adjacent
spans when the initial segmentation produces too many candidates relative to the
document length.  The guard targets a bounded number of chunks by packing
adjacent spans according to a word budget and, when configured, a soft target
band.  These operations preserve the original order and never discard text.

Finally, the method validates the resulting partition.  If a requested
segmentation strategy fails to produce a valid full-coverage partition, the
system falls back to fixed-token chunking using tokenizer offsets when
available, and whitespace token spans otherwise.  Thus every downstream scoring
step receives a well-formed chunk set $\mathcal{C}$ with stable chunk ids and
recoverable character spans.

\subsection{Submodular Function}

We define a set objective $F(S)$ over selected chunk subsets.  In the executable
method, the objective contains three terms: confidence, effectiveness, and
collaboration.  These terms correspond to three complementary requirements for
a useful attribution set: it should support the target prediction, cover
non-redundant evidence, and contain information that cannot be replaced by the
unselected complement.

\paragraph{Confidence Score.}
The confidence term measures whether the selected text alone supports the
target label.  For a subset $S$, we query the LLM on $x_S$ and use the target
label probability as the confidence score:
\[
    f_{\mathrm{conf}}(S)=p_{\theta}(y \mid x_S).
\]
This term rewards subsets that are sufficient for the model to assign high
probability to the target label.  It is target-conditioned: the same text subset
may receive different scores for different labels because the verbalizer
likelihood is computed with respect to the label being explained.

\paragraph{Effectiveness Score.}
Confidence alone can select repeated evidence.  For example, several chunks may
all express the same sentiment or repeat the same premise, producing high target
probability without adding new information.  The effectiveness term discourages
this redundancy by rewarding semantic diversity among the selected chunks.  Let
$\phi(c_i)$ be the normalized embedding of chunk $c_i$ and
$d(c_i,c_j)=1-\cos(\phi(c_i),\phi(c_j))$.  We define
\[
    f_{\mathrm{eff}}(S)=
    \begin{cases}
    0, & |S|\leq 1,\\
    \sum_{i\in S}\min_{j\in S, j\neq i} d(c_i,c_j), & |S|>1.
    \end{cases}
\]
Each selected chunk contributes its distance to the nearest other selected
chunk.  A redundant chunk has a nearby neighbor and therefore a small
contribution, while a complementary chunk increases the coverage of distinct
semantic content.  This term is computed from a precomputed chunk distance
matrix, so adding a candidate during search does not require another LLM
probability query.

\paragraph{Collaboration Score.}
Some chunks are important because they work together with other evidence, even
if their isolated confidence score is weak.  To capture such necessity, the
collaboration term evaluates how much the unselected complement loses the
semantic state of the original input after $S$ is removed.  Let
$\phi(x)$ be the embedding of the full text and $\phi(x_{\bar{S}})$ be the
embedding of the complement.  We define
\[
    f_{\mathrm{collab}}(S)=1-\cos(\phi(x_{\bar{S}}),\phi(x)).
\]
When $S=\emptyset$, the complement is the full text and the score is close to
zero.  As the search selects chunks whose removal changes the remaining text
substantially, the complement becomes less aligned with the original input and
the collaboration score increases.  This term therefore favors chunks that are
necessary to reconstruct the full input's semantic representation, not merely
chunks that are individually predictive.

\paragraph{Combined Objective.}
The final objective is a weighted sum of the three active scores:
\[
    F(S)=
    \lambda_{\mathrm{conf}} f_{\mathrm{conf}}(S)
    +\lambda_{\mathrm{eff}} f_{\mathrm{eff}}(S)
    +\lambda_{\mathrm{collab}} f_{\mathrm{collab}}(S).
\]
In our default active configuration, all three weights are set to one, giving
\[
    F(S)=
    f_{\mathrm{conf}}(S)
    +f_{\mathrm{eff}}(S)
    +f_{\mathrm{collab}}(S).
\]
The objective balances three complementary desiderata: selected chunks should
preserve the target prediction, cover diverse evidence rather than duplicates,
and remove information that the remaining complement cannot easily replace.

For efficiency, objective evaluations are cached by normalized subset keys.
When the search needs to evaluate all one-step augmentations
$S\cup\{i\}$, the implementation materializes the required subset texts,
complement texts, label probabilities, and embeddings in batches whenever
possible.  This preserves the exact objective while reducing repeated LLM
forward passes.

\subsection{Minimum Sufficient Explanation}

The explanation is obtained by approximately solving a cardinality-constrained
subset selection problem:
\[
    S^{\star}=\arg\max_{S\subseteq \{1,\ldots,n\},\ |S|\leq k} F(S).
\]
Here $k$ is a small explanation budget.  In practice, the method also returns a
complete chunk ranking: the first $k$ chunks form the selected explanation, and
the remaining chunks are appended according to their singleton gains so that
faithfulness metrics can evaluate arbitrary top-percent prefixes.

\paragraph{Forward Greedy Search.}
The standard search strategy starts from the empty set $S_0=\emptyset$.  At
step $t$, it scans every candidate chunk that has not yet been selected and
computes its conditional marginal gain
\[
    \Delta(i\mid S_{t-1})=
    F(S_{t-1}\cup\{i\})-F(S_{t-1}).
\]
The selected chunk is
\[
    i_t=\arg\max_{i\notin S_{t-1}}\Delta(i\mid S_{t-1}),
    \qquad
    S_t=S_{t-1}\cup\{i_t\}.
\]
Ties are resolved deterministically by choosing the smaller chunk id.  The
selected order $(i_1,\ldots,i_k)$ is the primary attribution ranking: an early
chunk is important not only because it scores well alone, but because it gives
the largest improvement under the evidence already selected.  Each search step
records the chosen chunk id, its marginal gain, the total objective value of the
augmented set, and the three score components of that augmented set.

\paragraph{Bidirectional Search.}
The bidirectional variant augments the forward greedy process with a reverse
pruning step.  It maintains two sets: the selected evidence set $S$ and a
removed set $R$ of chunks that are no longer considered candidates.  Each
iteration first performs the same forward addition as above over
\[
    \{1,\ldots,n\}\setminus(S\cup R).
\]
After adding the best chunk, the method evaluates which remaining chunk has the
smallest contribution to the current unremoved universe.  Let
$U_R=\{1,\ldots,n\}\setminus R$.  For every unselected and unremoved candidate
$j$, its reverse contribution is
\[
    \Gamma(j\mid U_R)=F(U_R)-F(U_R\setminus\{j\}).
\]
The algorithm removes the chunk with the smallest $\Gamma(j\mid U_R)$ from
future consideration:
\[
    j_t=\arg\min_{j\in U_R\setminus S}\Gamma(j\mid U_R),
    \qquad
    R \leftarrow R\cup\{j_t\}.
\]
The intuition is that a chunk whose deletion barely changes the objective is
unlikely to be necessary for a compact sufficient explanation.  By alternating
between adding the most useful chunk and discarding the least useful remaining
chunk, bidirectional search concentrates later scoring rounds on a smaller and
more relevant candidate pool.  As in forward greedy search, ties are
deterministic and only forward additions appear in the final explanation trace.

\paragraph{Explanation Construction.}
After search terminates, the selected text is formed by concatenating the chosen
chunks in their original document order, not in the order in which they were
selected.  This keeps the explanation readable and preserves local context.  The
importance ranking, however, follows the greedy selection order because that is
the order in which chunks deliver conditional marginal utility.  For chunks not
selected within the budget, we compute singleton gains
$F(\{i\})-F(\emptyset)$ and append them in descending order.  The result is a
complete ranking over all chunks, a selected subset of size at most $k$, and a
score trace that makes the contribution of each selected step explicit.
